\lecture{13}{Advanced Architectures in Relational Deep Learning}{Thursday, November 5, 2026}\label{ch:lec13}

% Local TikZ styles (prefix lec13)
\tikzset{
  lec13ent/.style={circle, draw=black!70, thick, fill=Salmon!22, minimum size=4.2mm,
                   inner sep=0pt},
  lec13jun/.style={circle, draw=black!70, thick, fill=Periwinkle!55, minimum size=4.2mm,
                   inner sep=0pt},
  lec13lbl/.style={font=\scriptsize\sffamily},
  lec13fk/.style={thick, black!40, -{Stealth[length=1.5mm]}},
  lec13cu/.style={circle, draw=black!70, thick, fill=ForestGreen!45, minimum size=5mm,
                  inner sep=0pt, font=\scriptsize},
  lec13tr/.style={rectangle, rounded corners=1pt, draw=black!70, thick, fill=RoyalBlue!30,
                  minimum size=4.6mm, inner sep=0pt, font=\scriptsize},
  lec13pr/.style={rectangle, draw=black!70, thick, fill=Goldenrod!55, minimum size=4.6mm,
                  inner sep=0pt, font=\scriptsize},
  lec13lnk/.style={thick, black!35, densely dotted, -{Stealth[length=1.5mm]}},
  lec13hop/.style={line width=1.6pt, black!35, -{Stealth[length=2.2mm]}},
  lec13hop2/.style={line width=1.6pt, black!70, -{Stealth[length=2.2mm]}},
  lec13cap/.style={font=\scriptsize\sffamily, align=center},
  lec13enc/.style={draw=black!80, fill=black!60, text=white, rounded corners=1.5pt,
                   font=\tiny\sffamily, minimum width=15mm, minimum height=5mm,
                   align=center, inner sep=1pt},
  lec13seed/.style={circle, fill=black!85, minimum size=2.2mm, inner sep=0pt},
  lec13t1/.style={circle, fill=Brown!80, minimum size=2.2mm, inner sep=0pt},
  lec13t2/.style={circle, fill=Gray!80, minimum size=2.2mm, inner sep=0pt},
  lec13up/.style={black!50, -{Stealth[length=1.3mm]}},
}

\textbf{Outline.} Recap of RDL (databases as graphs, training tables, why not LLMs, RDL
gains over LLM) $\to$ next-generation architectures (multi-modal, multi-scale databases;
structural patterns) $\to$ RelGNN: composite message passing along atomic routes $\to$
RelGT: Relational Graph Transformer (multi-element tokens) $\to$ relational foundation
model.

% ---------------------------------------------------------------------
\section{Recap: relational deep learning}\label{sec:lec13-recap}

\subsection{RDL recap}\label{sec:lec13-rdlrecap}
\slides{3--13}

All as in \cref{ch:lec12}:
\begin{itemize}
  \item Slide~4: data stored in relational databases (commerce, finance, social media
    schemas).
  \item Slide~5: predictions on relational data (Users--Sales--Products): which
    products will a user purchase in the next 7 days; will an active user churn in the
    next 90 days; total sales for each product in the next 30 days.
  \item Slides~6--7: insight: a database is a graph (schema $\to$ Users--Sales--Products
    graph).
  \item Slide~8: just do ML on a graph (node-, link-, graph-level tasks).
  \item Slide~9: graph ML problem solving pipeline (Relational DB $\to$ Graph Problem
    $\to$ Graph ML $\to$ Solutions).
  \item Slide~10: \textbf{training table} (entity ID, time, labels); classification,
    regression, multi-class.
  \item Slide~11: RDL pipeline (a) rel.\ tables with training table, (b) entities linked
    by foreign keys, (c) relational entity graph, (d) GNN.
  \item Slide~12: connect the training table: training labels with timestamps attached
    to the graph.
  \item Slide~13: GNN on the entity graph: a node's neighborhood defines a
    (time-dependent) computation graph.
\end{itemize}

\subsection{Why not use an LLM?}\label{sec:lec13-llm}
\slides{14--18}

\begin{itemize}
  \item Slides~14--17 as in \cref{ch:lec12}: some have tried (Wydmuch, Borchmann,
    Grali\'nski, Snowflake AI Research), the serialized prompt, and why not LLMs (not
    trained on the relational nature of databases; database prediction is \textbf{NOT}
    sequence modeling; \textbf{NOT} trained to predict the future; context size is
    limited).
  \item \textbf{New -- RDL gains over LLM} (slide~18): entity classification on the
    RelBench test set, ICL Llama 3.2 3B vs RelGNN (\cref{tab:lec13-llm}). RelGNN wins on
    all tasks except rel-f1. ``Stay tuned for more discussion on LLMs later!''
\end{itemize}

\begin{table}[htbp]
  \centering\small
  \caption{Entity classification results (ROC-AUC(\%), higher is better) on RelBench test
  set; best values in bold (slide~18).}
  \label{tab:lec13-llm}
  \begin{tabular}{@{}llccc@{}}
    \toprule
    Dataset & Task & ICL Llama 3.2 3B & RelGNN (ours) & Relative Gain\\
    \midrule
    \multirow{2}{*}{rel-amazon} & user-churn & 62.55 & \textbf{70.99} & 13\%\\
      & item-churn & 73.41 & \textbf{82.64} & 13\%\\
    \multirow{2}{*}{rel-avito} & user-visits & 53.36 & \textbf{66.18} & 24\%\\
      & user-clicks & 54.07 & \textbf{68.23} & 26\%\\
    \multirow{2}{*}{rel-event} & user-repeat & 70.11 & \textbf{79.61} & 14\%\\
      & user-ignore & 68.65 & \textbf{86.18} & 26\%\\
    \multirow{2}{*}{rel-f1} & driver-dnf & \textbf{80.03} & 75.29 & $-6$\%\\
      & driver-top3 & \textbf{87.11} & 85.69 & $-2$\%\\
    rel-hm & user-churn & 63.81 & \textbf{70.93} & 11\%\\
    \multirow{2}{*}{rel-stack} & user-engagement & 81.23 & \textbf{90.75} & 12\%\\
      & user-badge & 79.99 & \textbf{88.98} & 11\%\\
    rel-trial & study-outcome & 59.17 & \textbf{71.24} & 20\%\\
    \bottomrule
  \end{tabular}
\end{table}

% ---------------------------------------------------------------------
\section{Next-generation architectures for RDL}\label{sec:lec13-nextgen}

\subsection{The need for specialized models}\label{sec:lec13-need}
\slides{19--22}

\begin{itemize}
  \item Relational databases: \textbf{multi-modal}, \textbf{multi-scale} information.
  \item \textbf{Entity level:} text, image, numerical, categorical, timestamp. Schema:
    Products (ProductID, Description, Image, Size), Transactions (TransactionID,
    ProductID, \textbf{Timestamp}, CustomerID, Price), Customers (CustomerID, Name);
    edges ``Bought which'', ``Bought by''.
  \item \textbf{Database level:} \textbf{relational structure}, \textbf{temporal
    structure}.
\end{itemize}

\subsection{Relational DBs structure}\label{sec:lec13-patterns}
\slides{23}

\begin{itemize}
  \item \term{Tripartite patterns}: Users -- Sales -- Products (the graph of
    \cref{ch:lec12}).
  \item \term{Cycle patterns}: two users who both bought the shirt (\$20) and the shoe
    (\$80) close a cycle user -- sale -- shirt -- sale -- user -- sale -- shoe -- sale --
    user.
  \item \textbf{Structure changes over time!}
\end{itemize}

% ---------------------------------------------------------------------
\section{RelGNN: composite message passing for RDL}\label{sec:lec13-relgnn}

\subsection{Relational entity graph vs heterogeneous graph}\label{sec:lec13-junction}
\slides{24--28}

\begin{itemize}
  \item Same schema (Products, Transactions, Customers; ``Bought which'', ``Bought by'')
    as two graphs:
    \begin{itemize}
      \item \textbf{standard heterogeneous graph}: customer $\xrightarrow{\text{purchase}}$
        product;
      \item \textbf{relational entity graph}: customer $\leftarrow$ transaction $\to$
        product.
    \end{itemize}
  \item \textbf{Observation:} edges in relational entity graphs are \textbf{defined by
    primary--foreign key links} that merely record table connectivity \textbf{without
    semantics}.
  \item \term{Junction tables} are introduced, decomposing each \textbf{many-to-many
    association} into \textbf{a pair of one-to-many links}
    (\cref{fig:lec13-junction}).
\end{itemize}

\begin{figure}[htbp]
  \centering
  \begin{tikzpicture}[scale=0.85, transform shape]
    \begin{scope}
      \node[lec13cu, label={[lec13lbl]left:1}] (c1) at (0,1.2) {};
      \node[lec13cu, label={[lec13lbl]left:2}] (c2) at (0,0) {};
      \node[lec13pr, label={[lec13lbl]right:1}] (p1) at (2.2,1.2) {};
      \node[lec13pr, label={[lec13lbl]right:2}] (p2) at (2.2,0) {};
      \draw[thick, Lavender!80!black] (c1)--(p1) (c1)--(p2) (c2)--(p1);
      \node[lec13cap] at (0,-0.6) {customer};
      \node[lec13cap] at (2.2,-0.6) {product};
      \node[lec13cap] at (1.1,-1.2) {(a) many-to-many association};
    \end{scope}
    \begin{scope}[xshift=4.6cm]
      \node[lec13cu, label={[lec13lbl]left:1}] (d1) at (0,1.2) {};
      \node[lec13cu, label={[lec13lbl]left:2}] (d2) at (0,0) {};
      \node[lec13tr, label={[lec13lbl]above:1}] (t1) at (1.9,1.9) {};
      \node[lec13tr, label={[lec13lbl]above:2}] (t2) at (1.9,0.9) {};
      \node[lec13tr, label={[lec13lbl]above:3}] (t3) at (1.9,-0.1) {};
      \node[lec13pr, label={[lec13lbl]right:1}] (q1) at (3.8,1.2) {};
      \node[lec13pr, label={[lec13lbl]right:2}] (q2) at (3.8,0.3) {};
      \draw[lec13lnk] (t1)--(d1); \draw[lec13lnk] (t2)--(d1); \draw[lec13lnk] (t3)--(d2);
      \draw[lec13lnk] (t1)--(q2); \draw[lec13lnk] (t2)--(q2); \draw[lec13lnk] (t3)--(q1);
      \node[lec13cap] at (0,-0.6) {customer};
      \node[lec13cap] at (1.9,-0.6) {transaction};
      \node[lec13cap] at (3.8,-0.6) {product};
      \node[lec13cap] at (1.9,-1.2) {(b) junction table: pair of one-to-many links};
    \end{scope}
  \end{tikzpicture}
  \caption{A many-to-many customer--product association becomes, in the relational
  entity graph, a transaction junction table with two one-to-many links (slide~28).}
  \label{fig:lec13-junction}
\end{figure}

\subsection{Junction tables in schema graphs}\label{sec:lec13-schema}
\slides{29--33}

\begin{itemize}
  \item Slides~29--30: RelBench and its 7 diverse datasets, as in \cref{ch:lec12}.
  \item \textbf{rel-amazon}: customer -- \emph{review} -- product; \textbf{rel-hm}:
    customer -- \emph{transactions} -- article (purple = junction table).
  \item \textbf{rel-F1}: junction tables constructor\_standings, constructor\_results,
    standings, results, qualifying link constructors, races, drivers; circuits $\to$
    races (\cref{fig:lec13-f1}a).
  \item Two junction patterns:
    \begin{itemize}
      \item \term{bridge structure}: standings $\to$ races, drivers;
      \item \term{hub structure}: results $\to$ races, constructors, drivers.
    \end{itemize}
\end{itemize}

\begin{figure}[htbp]
  \centering
  \begin{tikzpicture}[scale=0.85, transform shape]
    \begin{scope}
      \node[lec13jun, label={[lec13lbl]above:constructor\_standings}] (cs) at (0.4,3.0) {};
      \node[lec13jun, label={[lec13lbl]above:constructor\_results}] (cr) at (3.2,3.0) {};
      \node[lec13jun, label={[lec13lbl]above:standings}] (st) at (5.8,3.0) {};
      \node[lec13ent, label={[lec13lbl]left:constructors}] (co) at (0,1.5) {};
      \node[lec13ent, label={[lec13lbl]right:races}] (ra) at (3.2,1.5) {};
      \node[lec13ent, label={[lec13lbl]right:drivers}] (dr) at (6.4,1.5) {};
      \node[lec13jun, label={[lec13lbl]below:results}] (re) at (1.8,0) {};
      \node[lec13ent, label={[lec13lbl]below:circuits}] (ci) at (3.2,0) {};
      \node[lec13jun, label={[lec13lbl]below:qualifying}] (qu) at (4.6,0) {};
      \foreach \a/\b in {cs/co, cs/ra, cr/co, cr/ra, st/ra, st/dr, re/co, re/ra, re/dr,
                         ci/ra, qu/co, qu/ra, qu/dr}
        \draw[lec13fk] (\a) -- (\b);
      \node[lec13cap] at (3.2,-1.05) {(a) primary--foreign key relation of rel-f1};
    \end{scope}
    \begin{scope}[xshift=9.4cm, yshift=0.2cm]
      \node[lec13ent, label={[lec13lbl]below:circuits}] (a1) at (0,2.9) {};
      \node[lec13ent, label={[lec13lbl]below:races}] (a2) at (2.4,2.9) {};
      \draw[lec13fk] (a1) -- (a2);
      \node[lec13ent, label={[lec13lbl]below:races}] (b1) at (0,1.6) {};
      \node[lec13jun, label={[lec13lbl]below:standings}] (b2) at (1.2,1.6) {};
      \node[lec13ent, label={[lec13lbl]below:drivers}] (b3) at (2.4,1.6) {};
      \draw[lec13fk] (b2) -- (b1); \draw[lec13fk] (b2) -- (b3);
      \node[lec13ent, label={[lec13lbl]below:races}] (c1) at (0,0.3) {};
      \node[lec13jun, label={[lec13lbl]below:results}] (c2) at (1.2,0.3) {};
      \node[lec13ent, label={[lec13lbl]below:drivers}] (c3) at (2.4,0.3) {};
      \draw[lec13fk] (c2) -- (c1); \draw[lec13fk] (c2) -- (c3);
      \node[lec13cap] at (1.2,-1.25) {(b) example atomic routes};
    \end{scope}
  \end{tikzpicture}
  \caption{rel-F1 schema graph (pink: entity tables, purple: junction tables) and example
  atomic routes (slides~32--33, 35).}
  \label{fig:lec13-f1}
\end{figure}

\subsection{Limitation of message passing on the entity graph}\label{sec:lec13-limit}
\slides{34}

\begin{itemize}
  \item \textbf{(1) \emph{Imbalance}:} \textbf{junction nodes} are 1-hop neighbors and
    \textbf{more informative nodes} become 2-hop (customer $\leftarrow$ transaction (1st
    hop neighbor) $\leftarrow$ product (2nd hop neighbor)).
  \item \textbf{(2) \emph{Redundancy}:} information from the source node is passed to the
    junction node in the first hop and routed back to itself in the second,
    \textbf{aggregated with the junction and 2-hop node information} (customer $\to$
    transaction in the 1st hop, transaction $\to$ customer in the 2nd hop).
\end{itemize}

\subsection{RelGNN: atomic routes}\label{sec:lec13-atomic}
\slides{35}

\begin{definition}[Atomic route, informal]\label{def:lec13-atomic}
  An \term{atomic route} is a \textbf{simple path between pink node-types} (non-junction
  tables) in the primary--foreign key graph.
\end{definition}
\begin{itemize}
  \item Atomic routes can be \textbf{derived automatically} without manual intervention
    or domain-expert knowledge.
  \item Examples on rel-f1 (\cref{fig:lec13-f1}b): circuits $\to$ races; races
    $\leftarrow$ standings $\to$ drivers; races $\leftarrow$ results $\to$ drivers.
  \item Chen, Kanatsoulis, Leskovec, \emph{RelGNN: Composite Message Passing for
    Relational Deep Learning}, ICML 2025.
\end{itemize}

\subsection{RelGNN: composite message passing}\label{sec:lec13-cmp}
\slides{36}

\begin{definition}[Composite message passing]\label{def:lec13-cmp}
  \term{Composite message passing} aggregates messages along atomic routes in a
  \textbf{single step}:
  \begin{equation}\label{eq:lec13-cmp}
    \vm^{(l+1)}_{(\mathtt{dst},\mathtt{mid},\mathtt{src})}
    =\mathrm{AGGR}\Big(\vh^{(l)}_{\mathtt{dst}},\
      \{\!\{\mathrm{FUSE}(\vh^{(l)}_{\mathtt{mid}},\vh^{(l)}_{\mathtt{src}})\}\!\}\Big)
  \end{equation}
\end{definition}
\textbf{Instantiation example} ($\vh^{(l)}_{\mathtt{fuse}}$ = FUSE output):
\begin{align}
  \mathrm{FUSE}(\vh^{(l)}_{\mathtt{mid}},\vh^{(l)}_{\mathtt{src}})
    &=\mW_1\vh^{(l)}_{\mathtt{mid}}+\mW_2\vh^{(l)}_{\mathtt{src}}
    \label{eq:lec13-fuse}\\
  \mathrm{AGGR}(\vh^{(l)}_{\mathtt{dst}},\{\!\{\vh^{(l)}_{\mathtt{fuse}}\}\!\})
    &=\mW_{\mathrm{proj}}\vh^{(l)}_{\mathtt{dst}}
      +\sum_{\mathtt{fuse}\in\mathcal{N}(\mathtt{dst})}
       \alpha_{\mathtt{dst},\mathtt{fuse}}\mW_V\vh^{(l)}_{\mathtt{fuse}}
    \label{eq:lec13-aggr}
\end{align}
\begin{itemize}
  \item \textbf{Standard} vs \textbf{composite message passing} in
    \cref{fig:lec13-cmp}: transaction and product are FUSEd and reach the customer in a
    single hop.
\end{itemize}

\begin{figure}[htbp]
  \centering
  \begin{tikzpicture}[scale=0.85, transform shape]
    \begin{scope}
      \node[lec13cu] (c) at (0,0) {};
      \node[lec13tr] (t) at (2.2,0) {};
      \node[lec13pr] (p) at (4.4,0) {};
      \draw[lec13hop] (0.4,0.25) -- node[above, lec13lbl] {1st hop} (1.8,0.25);
      \draw[lec13hop2] (1.8,-0.25) -- node[below, lec13lbl] {2nd hop} (0.4,-0.25);
      \draw[lec13hop] (4.0,0) -- node[above, lec13lbl] {1st hop} (2.6,0);
      \node[lec13cap] at (0,-0.9) {customer};
      \node[lec13cap] at (2.2,-0.9) {transaction};
      \node[lec13cap] at (4.4,-0.9) {product};
      \node[lec13cap] at (2.2,-1.5) {(a) standard message passing};
    \end{scope}
    \begin{scope}[xshift=6.6cm]
      \node[lec13cu] (c) at (0,0) {};
      \node[lec13tr] (t) at (2.3,0) {};
      \node[lec13pr] (p) at (3.9,0) {};
      \draw[thick, rounded corners=2pt, black!45] (1.7,-0.45) rectangle (4.5,0.75);
      \node[lec13lbl, font=\scriptsize\sffamily\bfseries] at (3.1,0.5) {FUSE};
      \draw[lec13hop] (1.6,0) -- node[above, lec13lbl] {single hop} (0.4,0);
      \node[lec13cap] at (0,-0.9) {customer};
      \node[lec13cap] at (2.3,-0.9) {transaction};
      \node[lec13cap] at (3.9,-0.9) {product};
      \node[lec13cap] at (2.2,-1.5) {(b) composite message passing};
    \end{scope}
  \end{tikzpicture}
  \caption{Standard message passing on customer--transaction--product (imbalance and
  redundancy, slide~34) vs composite message passing along the atomic route
  (slide~36).}
  \label{fig:lec13-cmp}
\end{figure}

\subsection{RelGNN results}\label{sec:lec13-relgnn-res}
\slides{37}

Hetero-GNN vs RelGNN (bar heights read off the plots, approximate):
\begin{itemize}
  \item \textbf{rel-F1} -- predict if the driver will qualify in the top-3 for a race in
    the next 1 month; ROC-AUC(\%) $\approx75$ vs $\approx86$: \textbf{13\% improvement}.
  \item \textbf{rel-trial} -- predict the success rate of a trial site in the next 1
    year; MAE $\approx0.40$ vs $\approx0.30$: \textbf{25\% improvement}.
  \item \textbf{rel-amazon} -- predict the list of distinct items each customer will
    purchase and give a detailed review in the next 3 months; MAP(\%) $\approx0.47$ vs
    $\approx0.52$: \textbf{11\% improvement}.
\end{itemize}
\begin{keyidea}
  RelGNN achieves \textbf{SOTA performance} on the vast majority of tasks from RelBench,
  with improvements of \textbf{up to 25\%}.
\end{keyidea}

% ---------------------------------------------------------------------
\section{RelGT: Relational Graph Transformer}\label{sec:lec13-relgt}

\subsection{Recap: graph transformers}\label{sec:lec13-gtrecap}
\slides{38--44}

Header citation: You, Ying, Leskovec, \emph{Design Space of Graph Neural Networks},
NeurIPS 2020 (\cref{ch:lec04}). Content as in \cref{ch:lec08}:
\begin{itemize}
  \item Slides~39--40: key components of a Transformer: (1) tokenizing, (2) positional
    encoding, (3) self-attention. Key question: what should these be for a graph input?
  \item Slides~41--43: a graph Transformer must take (1) node features, (2) adjacency
    information, (3) edge features; many ways, different ``matchings'': node features
    $\leftrightarrow$ tokenizing, adjacency $\leftrightarrow$ positional encoding, edge
    features $\leftrightarrow$ self-attention.
  \item Slide~44: PEARL + Transformer in practice: sample node ids $\to$ GNN on each
    sample set $\to$ empirical expectation $\to$ concatenate with node features $\mX$
    $\to$ main GNN/Transformer $\to$ train PEARL + prediction model jointly.
\end{itemize}

\subsection{Observation: what a token represents}\label{sec:lec13-tokens}
\slides{45--48}

\begin{itemize}
  \item \textbf{(Text) Transformers:} an input token is represented with \textbf{2
    elements}: token id (\emph{feature}), position (\emph{structure}) -- input embedding
    + positional encoding.
  \item \textbf{Graph Transformers:} an input token is represented with \textbf{2
    elements}: node feature (\emph{feature}), graph positional encoding
    (\emph{structure}) (graph transformer layer, $\times H$ heads, $\times L$ layers).
    \begin{itemize}
      \item Most GTs are designed for \textbf{homogeneous, static graphs};
      \item for additional RDL complexity such as \textbf{heterogeneity} and
        \textbf{temporality}, this design needs to be improved.
    \end{itemize}
  \item \textbf{Graph Transformer inputs for RDL} (slide~48): (1) node features
    $\leftrightarrow$ tokenizing, (2) adjacency information $\leftrightarrow$ positional
    encoding; (3) edge features is struck out and replaced by \textbf{(3) node types}
    and \textbf{(4) timestamp}.
\end{itemize}

\subsection{Relational Graph Transformer: overview and challenges}\label{sec:lec13-relgt-ov}
\slides{49--53}

\begin{itemize}
  \item \textbf{Pipeline} (slide~49): relational entity graph $\to$ \textbf{tokens}
    (seed node + local context) $\to$ \textbf{transformer network} ($\times L$ layers;
    the seed node attends to \textbf{local tokens} and to \textbf{global centroids})
    $\to$ prediction (downstream layers).
  \item In the \term{Relational Graph Transformer} we use \textbf{multiple elements to
    represent the graph structure}: node feature, $\langle$multiple information
    elements$\rangle$ (= structure).
  \item \textbf{Limitations of existing PEs for RDL:} designed for homogeneous graphs; no
    temporal dynamics; computational overhead; graph changes over time.
  \item \textbf{Unique properties of relational graphs:} the relational entity graph is
    \textbf{large-scale}, \textbf{heterogeneous}, \textbf{temporal}.
    \begin{itemize}
      \item add a \textbf{categorical encoding} to capture the heterogeneity;
      \item add a \textbf{time encoding} to capture the temporal structure.
    \end{itemize}
  \item \textbf{How to encode structure.} Challenge: the relational entity graph is
    large-scale. Solution: \textbf{sample a subgraph around the seed node}.
  \item \term{Training (seed) nodes} correspond to the node-types with respect to which
    tasks are defined; we sample a \textbf{context (subgraph)} around each seed node.
\end{itemize}

\subsection{Token preparation (sampling stage)}\label{sec:lec13-sampling}
\slides{54--56}

\begin{itemize}
  \item For each training node, a \textbf{fixed set of $K$ tokens} from the local
    neighborhood (e.g., up to 2 hops) is selected through a \textbf{temporal-aware
    sampling} $\Rightarrow$ relational entity graph $\xrightarrow{\text{tokenization}}$
    $B\times K$ tokens for $B$ training nodes.
  \item Structure is encoded with a \textbf{GNN positional encoding} for each node in the
    context + \textbf{hop distance from seed node}.
  \item Each node in the $K$ set is represented by a \textbf{5-tuple}:
    \textbf{(node feature, node type, hop distance, time, GNN PE)}.
  \item Slide~56: seed nodes $\to$ local tokens (one row of tokens per seed node).
\end{itemize}

\subsection{Token feature encoders}\label{sec:lec13-encoders}
\slides{57--59}

\begin{itemize}
  \item Each element of the 5-tuple is passed to a \textbf{feature encoder} and the
    combination becomes the token features for the Transformer network
    (\cref{fig:lec13-enc}):
    \begin{itemize}
      \item node feature $\to$ Torch Frame Encoder; node type $\to$ nn Embedding; hop
        distance $\to$ nn Embedding; time $\to$ Time Encoder; subgraph $\to$ GNN
        Encoder;
      \item the element-wise encoder outputs are combined ($\oplus$) $\to$ tokens fed to
        the Transformer.
    \end{itemize}
  \item Slide~59: the \textbf{subgraph} input of the GNN Encoder is the sampled local
    subgraph around the seed node.
\end{itemize}

\begin{figure}[htbp]
  \centering
  \begin{tikzpicture}[scale=0.85, transform shape]
    % seed nodes and local tokens
    \foreach \y in {0,-0.38,-0.76} \node[lec13seed] at (-1.3,\y) {};
    \node[lec13cap, anchor=east] at (-1.55,-0.38) {seed nodes};
    \draw[-{Stealth[length=1.5mm]}, Magenta] (-0.95,-0.38) -- (-0.25,-0.38);
    \node[lec13seed] (s0) at (0,0) {};
    \node[lec13seed] at (0,-0.38) {}; \node[lec13seed] at (0,-0.76) {};
    \node[lec13t1] at (0.35,0) {}; \node[lec13t2] at (0.7,0) {};
    \node[font=\tiny] at (1.1,0) {$\cdots$};
    \node[lec13t1] at (1.5,0) {}; \node[lec13t2] at (1.85,0) {}; \node[lec13t2] at (2.2,0) {};
    \node[lec13cap, anchor=west] at (0.6,-0.45) {local tokens};
    % element-wise encoders
    \draw[thick, rounded corners=4pt, black!60] (-0.9,0.45) rectangle (9.3,2.75);
    \foreach \i/\inp/\enc in {0/node feature/Torch Frame\\Encoder,
                              1/node type/nn Embedding,
                              2/hop distance/nn Embedding,
                              3/time/Time Encoder,
                              4/subgraph/GNN Encoder} {
      \node[lec13lbl] (in\i) at ({0.05+2.05*\i},0.8) {\inp};
      \node[lec13enc] (e\i) at ({0.05+2.05*\i},1.65) {\enc};
      \draw[lec13up] (in\i) -- (e\i);
      \draw[lec13up] (e\i) -- ({0.05+2.05*\i},2.45);
    }
    \draw[lec13up] (s0) -- (in0);
    \node[lec13cap, anchor=east, align=right] at (-1.0,1.65) {element-wise\\encoders};
    \node[circle, draw, thick, inner sep=0pt, minimum size=3mm] (sum) at (4.15,3.15) {};
    \draw (sum.north) -- (sum.south) (sum.west) -- (sum.east);
    \draw[lec13up] (4.15,2.75) -- (sum);
    % output tokens
    \node[lec13seed] (o0) at (3.1,3.75) {};
    \node[lec13t1] at (3.45,3.75) {}; \node[lec13t2] at (3.8,3.75) {};
    \node[font=\tiny] at (4.2,3.75) {$\cdots$};
    \node[lec13t1] at (4.6,3.75) {}; \node[lec13t2] at (4.95,3.75) {}; \node[lec13t2] at (5.3,3.75) {};
    \draw[lec13up] (sum) -- (4.15,3.6);
    \draw[-{Stealth[length=1.5mm]}, Magenta] (5.6,3.75) -- (6.3,3.75)
      node[right, lec13cap, text=black] {tokens fed to Transformer};
  \end{tikzpicture}
  \caption{RelGT token feature encoders: each element of the 5-tuple goes through its own
  encoder and the outputs are combined into the token fed to the Transformer
  (slides~57--59).}
  \label{fig:lec13-enc}
\end{figure}

\subsection{Ablation studies and results}\label{sec:lec13-relgt-res}
\slides{60--61}

\begin{itemize}
  \item \textbf{Ablation} (\cref{tab:lec13-ablation}): relative drop (\%) in performance
    in RelGT after removing a model component; negative scores suggest the component is
    critical in RelGT, and vice-versa.
\end{itemize}
\begin{keyidea}[Takeaway, slide 60]
  Contribution of the \textbf{multi-element tokenization strategy} to capture arbitrary
  structure in relational data.
\end{keyidea}

\begin{table}[htbp]
  \centering\small
  \caption{Relative drop (\%) in performance in RelGT after removing a model component
  (slide~60; full results in the paper's Table~7).}
  \label{tab:lec13-ablation}
  \begin{tabular}{@{}llrrrrr@{}}
    \toprule
    Dataset & Task & No Global & No GNN & No Node & No Hop & No Relative\\
     & & Module & PE & Type & Distance & Time\\
    \midrule
    rel-avito & ad-ctr & $-6.00$ & $-1.14$ & $-7.14$ & $-3.43$ & $-9.14$\\
    rel-avito & user-clicks & 7.79 & $-15.19$ & 4.96 & 5.72 & 8.32\\
    rel-avito & user-visits & $-0.32$ & $-2.35$ & $-0.08$ & 0.42 & $-0.72$\\
    rel-event & user-ignore & $-1.28$ & 0.15 & $-0.09$ & 0.69 & $-0.06$\\
    rel-trial & study-outcome & $-2.08$ & $-1.66$ & 3.80 & $-0.37$ & 2.54\\
    rel-trial & site-success & $-19.08$ & $-9.23$ & $-2.94$ & $-21.56$ & $-0.77$\\
    rel-amazon & user-churn & $-0.64$ & $-0.78$ & 0.16 & 0.06 & $-2.20$\\
    rel-hm & item-sales & $-9.33$ & $-17.35$ & $-12.69$ & 0.93 & $-77.24$\\
    \midrule
     & \textbf{Average} & $-3.87$ & $-5.95$ & $-1.75$ & $-2.19$ & $-9.91$\\
    \bottomrule
  \end{tabular}
\end{table}

\begin{itemize}
  \item \textbf{Results} (\cref{tab:lec13-relgt}): test set results on the entity
    regression and classification tasks in RelBench. RDL: HeteroGNN baseline; HGT:
    Heterogeneous GT; PE: Laplacian Positional Encodings; \% Rel.\ Gain = percentage
    improvement over the RDL baseline.
\end{itemize}
\begin{keyidea}[Takeaway, slide 61]
  Relational Graph Transformers \textbf{improve over GNNs}.
\end{keyidea}

\begin{table}[tb]
  \centering\footnotesize
  \caption{RelGT results on RelBench, best values in bold as on slide~61. Top: (a) MAE for
  entity regression, lower is better. Bottom: (b) AUC for entity classification, higher
  is better.}
  \label{tab:lec13-relgt}
  \begin{tabular}{@{}llccccr@{}}
    \toprule
    Dataset & Task & RDL & HGT & HGT+PE & RelGT (ours) & \% Rel.\ Gain\\
    \midrule
    \multicolumn{7}{@{}l}{\emph{(a) MAE for entity regression (lower is better)}}\\
    rel-f1 & driver-position & 4.022 & 4.1598 & 4.2358 & \textbf{3.9170} & 2.61\\
    rel-avito & ad-ctr & 0.041 & 0.0441 & 0.0494 & \textbf{0.0345} & 15.85\\
    rel-event & user-attendance & 0.258 & 0.2635 & 0.2562 & \textbf{0.2502} & 2.79\\
    \multirow{2}{*}{rel-trial} & study-adverse & 44.473 & 43.3253 & \textbf{42.4622}
      & 43.9923 & 1.08\\
      & site-success & 0.400 & 0.4374 & 0.4431 & \textbf{0.3263} & 18.43\\
    \multirow{2}{*}{rel-amazon} & user-ltv & 14.313 & 15.3804 & 15.9296
      & \textbf{14.2665} & 0.32\\
      & item-ltv & 50.053 & 56.1384 & 55.6211 & \textbf{48.9222} & 2.26\\
    rel-stack & post-votes & \textbf{0.065} & 0.0679 & 0.0680 & \textbf{0.0654} & $-0.62$\\
    rel-hm & item-sales & 0.056 & 0.0655 & 0.0641 & \textbf{0.0536} & 4.29\\
    \midrule
    \multicolumn{7}{@{}l}{\emph{(b) AUC for entity classification (higher is better)}}\\
    \multirow{2}{*}{rel-f1} & driver-dnf & 0.7262 & 0.7142 & 0.7109 & \textbf{0.7587}
      & 4.48\\
      & driver-top3 & 0.7554 & 0.6389 & 0.8340 & \textbf{0.8352} & 10.56\\
    \multirow{2}{*}{rel-avito} & user-clicks & 0.6590 & 0.6584 & 0.6387
      & \textbf{0.6830} & 3.64\\
      & user-visits & 0.6620 & 0.6426 & 0.6507 & \textbf{0.6678} & 0.88\\
    \multirow{2}{*}{rel-event} & user-repeat & \textbf{0.7689} & 0.6717 & 0.6590
      & 0.7609 & $-1.04$\\
      & user-ignore & 0.8162 & \textbf{0.8348} & 0.8161 & 0.8157 & $-0.06$\\
    rel-trial & study-outcome & 0.6860 & 0.5679 & 0.5691 & \textbf{0.6861} & 0.01\\
    \multirow{2}{*}{rel-amazon} & user-churn & \textbf{0.7042} & 0.6608 & 0.6589
      & 0.7039 & $-0.04$\\
      & item-churn & \textbf{0.8281} & 0.7824 & 0.7840 & 0.8255 & $-0.31$\\
    \multirow{2}{*}{rel-stack} & user-engagement & 0.9021 & 0.8898 & 0.8852
      & \textbf{0.9053} & 0.35\\
      & user-badge & \textbf{0.8986} & 0.8652 & 0.8518 & 0.8632 & $-3.94$\\
    rel-hm & user-churn & \textbf{0.6988} & 0.6773 & 0.6491 & 0.6927 & $-0.87$\\
    \bottomrule
  \end{tabular}
\end{table}

% ---------------------------------------------------------------------
\section{Relational foundation model}\label{sec:lec13-rfm}

\subsection{Limitations of previous models}\label{sec:lec13-prev}
\slides{62--64}

\begin{itemize}
  \item Previous RDL models (labels attached to the Users--Sales--Products graph):
    \begin{itemize}
      \item they are \textbf{schema- and task-specific};
      \item they do not transfer knowledge from one database to another or one task to
        another;
      \item they do not generalize to new databases or tasks.
    \end{itemize}
  \item \textbf{Can we do better? Goal:} bring the success of \textbf{foundation models}
    (LLMs, ViT) to relational data (ViT: patch + position embedding $\to$ Transformer
    encoder $\to$ MLP head $\to$ class).
\end{itemize}

\subsection{Relational foundation model and RelGT as its backbone}\label{sec:lec13-kumorfm}
\slides{65--67}

\begin{itemize}
  \item \textbf{Key capabilities of Relational Foundation Models (RFMs)} (slide~65):
    \begin{itemize}
      \item applied to \textbf{new/unseen databases and schemas} with highly varying
        structural characteristics: (a) e-commerce (Users, Orders, Likes, Items), (b)
        finance (Customers, Transactions, Account, Branch), (c) health care (Patient,
        Record, Visit, Doctor);
      \item applied to \textbf{any predictive task type}, from one-off assessments
        (e.g.\ entity-level fraud prediction) to temporal predictive queries (e.g.\
        temporal recommendation prediction); example queries: ``What will user 123 buy
        next?'', ``Will user 321 churn soon?'', ``Is transaction 567 fraudulent?'',
        ``Will clinic 987 have high appointment demand?'', ``Is patient 234 at risk of
        hospital readmission?'';
      \item \textbf{generalize to new predictive tasks} and give accurate predictions
        without any task-specific model tuning;
      \item support prediction outputs, \textbf{explanations} into the reasoning
        process, and build trust through extensive quantitative \textbf{evaluation}.
    \end{itemize}
  \item \textbf{RelGT as the backbone for RFM} (slide~66; \emph{Introducing KumoRFM: A
    Foundation Model for In-Context Learning on Relational Data}, Kumo, May 20, 2025):
    \begin{itemize}
      \item query in PQL, e.g.\ \texttt{PREDICT COUNT(orders.*, 0, 7) FOR users.user\_id
        = 0} $\to$ PQL Parser $\to$ PQL Validator;
      \item from the relational database, a \textbf{backward-looking graph sampler}
        (input) and a \textbf{forward-looking graph sampler} (output) build the context
        and the prediction query:
        \begin{equation}\label{eq:lec13-icl}
          \text{context: }\big\{\calG^{\le\hat t}[\hat e]\big\}_{\hat t<t,\ \hat e\sim\calG},
          \ \big\{y^{\hat t}_{\hat e}\big\}_{\hat t<t,\ \hat e\sim\calG};
          \qquad
          \text{predict: }\calG^{\le t}[e]\ \to\ ?
        \end{equation}
      \item Relational Graph Transformer ($L$ times) $\to$ \textbf{In-Context Learning}
        (with ``Explain'').
    \end{itemize}
\end{itemize}
\begin{keyidea}
  \textbf{Relational Graph Transformer} is the backbone for the first Relational
  Foundation Model with \textbf{zero-shot capabilities powered by ICL}.
\end{keyidea}
\begin{itemize}
  \item \textbf{Time to first prediction} (slide~67) on node-level classification tasks
    (user-engagement, user-badge, user-churn (amazon), item-churn, driver-dnf,
    driver-top3, user-churn (hm), study-outcome), Data Scientist vs RDL vs KumoRFM:
    (a) hours of human labor (in seconds, log-scaled), (b) lines of code (log-scaled).
    KumoRFM is \textbf{orders of magnitude faster} ($\approx1$ second vs $\approx30$
    minutes vs $\approx12.3$ hours) and requires \textbf{zero-code} to get to accurate
    predictions (1 LoC vs $\approx56$ LoC vs $\approx878$ LoC).
\end{itemize}
\begin{caution}
  Slide~67's caption (``respectively'') pairs $\approx$30~min\,/\,$\approx$56~LoC with Data
  Scientist and $\approx$12.3~h\,/\,$\approx$878~LoC with RDL; the bars show the reverse
  (Data Scientist bars longest).
\end{caution}

% ---------------------------------------------------------------------
\bigskip\phantomsection
\addcontentsline{toc}{section}{Lecture 13 at a glance}
\begin{summarybox}[Lecture 13 at a glance]
  \begin{itemize}
    \item \textbf{RDL vs LLM:} RelGNN beats ICL Llama 3.2 3B on RelBench entity
      classification (gains 11--26\%), except rel-f1 ($-6$\%, $-2$\%).
    \item \textbf{Relational DBs:} multi-modal (text, image, numerical, categorical,
      timestamp), multi-scale (relational + temporal structure); tripartite and cycle
      patterns; structure changes over time.
    \item \textbf{Entity-graph edges} = PK--FK links without semantics; \textbf{junction
      tables} split many-to-many into two one-to-many links (bridge / hub) $\Rightarrow$
      \textbf{imbalance} (informative nodes at 2 hops) and \textbf{redundancy} (own
      information routed back).
    \item \textbf{RelGNN:} atomic route = simple path between non-junction node types;
      composite message passing
      $\vm^{(l+1)}_{(\mathtt{dst},\mathtt{mid},\mathtt{src})}
      =\mathrm{AGGR}(\vh^{(l)}_{\mathtt{dst}},\{\!\{\mathrm{FUSE}(\vh^{(l)}_{\mathtt{mid}},
      \vh^{(l)}_{\mathtt{src}})\}\!\})$, FUSE
      $=\mW_1\vh^{(l)}_{\mathtt{mid}}+\mW_2\vh^{(l)}_{\mathtt{src}}$;
      SOTA on most RelBench tasks, up to 25\%.
    \item \textbf{RelGT:} GTs are built for homogeneous static graphs; RelGT samples
      $K$ tokens per seed node (temporal-aware, e.g.\ up to 2 hops), each a 5-tuple (node
      feature, node type, hop distance, time, GNN PE) with its own encoder, combined;
      attention to local tokens + global centroids. Ablation: removing any component
      lowers performance on average ($-9.91$\% without relative time); improves over
      GNNs.
    \item \textbf{RFM:} previous models are schema- and task-specific; KumoRFM uses RelGT
      as backbone with in-context learning (zero-shot): $\approx1$ s and 1 LoC to first
      prediction.
  \end{itemize}
\end{summarybox}
